\section{序論}
11月4日07時から5日07時（日本時間）のXでは、中国発の大規模モデルと、AIを実際の作業へ組み込む技術が目立った。TencentのHy4は視覚理解とコンピュータ操作を前面に掲げ、モデルの規模だけでなく、作業環境を扱える能力が競争軸になっていることを示した。

同時に、AnthropicはClaudeを使った構造生物学の研究事例を紹介し、MetaによるConductor買収は複数のコーディングエージェントを束ねる技術への関心を映した。Mistralのセキュリティ対応デモも、エージェントが専門作業をどこまで担えるかという問いにつながる。

本号はGrokの07時収集結果5件を記載順に記事化した。性能や実用性については、投稿・発表側の主張と第三者による検証を区別して読む必要がある。
